% 子文件：5.复变函数的积分.tex
% 对应教材：《复变函数与积分变换》(哈工大, 第三版) 第 3 章
% 对应讲义：包益欣《复变函数》第 3 章（复积分的定义 / 柯西积分定理及应用 / 柯西积分公式及证明 / 其应用）
% 节次：§3.1 复变函数积分的概念（定义 / 性质 / 存在条件与计算）
%       §3.2 柯西积分定理（3.2.1 定理 / 3.2.2 不定积分 / 3.2.3 复合闭路定理）
%       §3.3 柯西积分公式（3.3.1 公式 / 3.3.2 高阶导数公式 / 3.3.3 几个重要的推论）
%       附加：§3.4 解析函数模的性质（最大模、最小模原理，据讲义 §1.4）
% 由 1.main.tex 用 \input 引入

\chapter{复变函数的积分}

\section{复变函数积分的概念}

\subsection{积分的定义}
先约定曲线。闭区间 \([a,b]\) 上的连续映射
\[
C:\ z=z(t)=x(t)+\mathrm{i}y(t),\qquad a\leqslant t\leqslant b
\]
称为复平面上由 \(z(a)\) 到 \(z(b)\) 的\textbf{有向曲线}，\(z(a)\)、\(z(b)\) 为起点与终点。称 \(z=z(a+b-t)\ (a\leqslant t\leqslant b)\) 为 \(C\) 的\textbf{反向曲线}，记 \(C^{-}\)。若 \(z'(t)\) 连续且处处不为零，称 \(C\) \textbf{光滑}；有限条光滑曲线首尾相接称为\textbf{逐段光滑}。若 \(z(a)=z(b)\) 且 \(t\in(a,b)\) 时曲线无重点，称 \(C\) 为\textbf{简单闭曲线}；不作说明时规定其\textbf{逆时针}方向为正向。

设 \(C\) 是逐段光滑有向曲线，\(f(z)\) 在 \(C\) 上有定义。在 \([a,b]\) 上插入分点
\[
a=t_{0}<t_{1}<\cdots<t_{n}=b,
\]
记 \(z_{k}=z(t_{k})\)，\(\Delta z_{k}=z_{k}-z_{k-1}\)，在每个小弧 \(z_{k-1}z_{k}\) 上任取一点 \(\zeta_{k}\)，作\textbf{黎曼和}
\[
\sum_{k=1}^{n}f(\zeta_{k})\,\Delta z_{k}.
\]
令 \(\lambda=\max\limits_{1\leqslant k\leqslant n}|\Delta z_{k}|\)，若 \(\lambda\to 0\) 时该和式的极限存在且与分法与 \(\zeta_{k}\) 的取法无关，则称此极限为 \(f(z)\) 沿 \(C\) 的\textbf{复积分}，记
\[
\int_{C}f(z)\,\mathrm{d}z=\lim_{\lambda\to 0}\sum_{k=1}^{n}f(\zeta_{k})\,\Delta z_{k},
\]
\(C\) 为积分路径，\(f\) 为被积函数。沿闭曲线的积分记作 \(\oint_{C}f(z)\,\mathrm{d}z\)。

\subsection{积分的性质}
设所涉积分都存在，则
\begin{enumerate}[leftmargin=2em]
  \item \textbf{线性性}：\(\displaystyle\int_{C}\bigl[\alpha f(z)+\beta g(z)\bigr]\mathrm{d}z=\alpha\int_{C}f\,\mathrm{d}z+\beta\int_{C}g\,\mathrm{d}z\)（\(\alpha,\beta\) 为复常数）。
  \item \textbf{反向性}：\(\displaystyle\int_{C^{-}}f(z)\,\mathrm{d}z=-\int_{C}f(z)\,\mathrm{d}z\)。
  \item \textbf{路径可加性}：若 \(C\) 由 \(C_{1}\) 与 \(C_{2}\) 首尾相接而成，则 \(\displaystyle\int_{C}f\,\mathrm{d}z=\int_{C_{1}}f\,\mathrm{d}z+\int_{C_{2}}f\,\mathrm{d}z\)。
  \item \textbf{估值不等式}：设 \(|f(z)|\leqslant M\ (z\in C)\)，\(L\) 为 \(C\) 的长度，则
  \[
  \left|\int_{C}f(z)\,\mathrm{d}z\right|\leqslant\int_{C}|f(z)|\,|\mathrm{d}z|=\int_{C}|f(z)|\,\mathrm{d}s\leqslant ML .
  \]
\end{enumerate}

\corepoint{估值不等式是复积分里\emph{唯一}的“放缩工具”，凡证极限为零、证积分有界、证刘维尔定理都靠它。使用时两个量都要算准：\(M\) 必须是 \(|f|\) 在 \(C\) 上的\textbf{上界}（常取最大值），\(L\) 是路径长度（圆周为 \(2\pi R\)）。注意被积函数在路径上含参数时 \(M\) 也随参数变，要一并估计。}

\subsection{积分的存在条件与计算}
\textbf{定理}\quad 若 \(f(z)\) 在逐段光滑曲线 \(C\) 上连续，则 \(\displaystyle\int_{C}f(z)\,\mathrm{d}z\) 存在。

\textbf{计算方法一（参数化，最基本）}\quad 设 \(C:z=z(t)\ (a\leqslant t\leqslant b)\) 光滑，则
\keyformula{\int_{C}f(z)\,\mathrm{d}z=\int_{a}^{b}f\bigl[z(t)\bigr]\,z'(t)\,\mathrm{d}t}
逐段光滑时按段积分再相加。这一步把复积分化为定积分（被积函数为复值函数，实部虚部分别积分）。

\textbf{计算方法二（化为两个第二类线积分）}\quad 设 \(f=u+\mathrm{i}v\)，\(z=x+\mathrm{i}y\)，则
\keyformula{\int_{C}f(z)\,\mathrm{d}z=\int_{C}(u\,\mathrm{d}x-v\,\mathrm{d}y)+\mathrm{i}\int_{C}(v\,\mathrm{d}x+u\,\mathrm{d}y)}
它说明复积分存在 \(\iff\) 右端两个线积分都存在；也说明复积分与路径方向有关。

\textbf{两个基本结果}（后面反复使用）
\keyformula{\int_{C}\mathrm{d}z=z(b)-z(a);\qquad
\oint_{|z-z_{0}|=r}(z-z_{0})^{n}\,\mathrm{d}z=
\begin{cases}2\pi\mathrm{i}, & n=-1,\\ 0, & n\neq -1,\ n\in\mathbb{Z}\end{cases}}
前者由参数化直接得到（\(\int_{a}^{b}z'(t)\mathrm{d}t=z(b)-z(a)\)），后者把 \(z=z_{0}+r\mathrm{e}^{\mathrm{i}\theta}\) 代入即可验证。它是“绕奇点积分不为零”的根源。

\section{柯西积分定理}

\subsection{柯西积分定理}
\textbf{定理（柯西积分定理）}\quad 设 \(f(z)\) 在\textbf{单连通区域} \(D\) 内解析，\(C\) 为 \(D\) 内任一条逐段光滑闭曲线，则
\keyformula{\oint_{C}f(z)\,\mathrm{d}z=0}
常用的等价叙述：若 \(f\) 在简单闭曲线 \(C\) 上及其内部解析，则沿 \(C\) 的积分为零。

\textbf{条件不能削弱}：只要求 \(f\) 连续不够，必须解析；区域必须是单连通的，否则绕“洞”的积分一般不为零——在去心圆盘 \(0<|z|<1\) 内 \(f(z)=1/z\) 解析，但 \(\oint_{|z|=1}\mathrm{d}z/z=2\pi\mathrm{i}\neq 0\)。

\textbf{推论（积分与路径无关）}\quad 若 \(f\) 在单连通区域 \(D\) 内解析，则 \(\int_{C}f\,\mathrm{d}z\) 只与起点、终点有关，与路径无关。

\textbf{闭路变形原理}\quad 在积分值不变的条件下，闭曲线可以在区域内部连续变形，只要变形过程中\textbf{不越过奇点}。它是复合闭路定理的几何表述，也是“挖小圆”技巧的依据。

\subsection{不定积分}
\textbf{定理（原函数存在定理）}\quad 设 \(f(z)\) 在单连通区域 \(D\) 内解析，固定 \(z_{0}\in D\)，令
\[
F(z)=\int_{z_{0}}^{z}f(\xi)\,\mathrm{d}\xi
\]
（积分沿 \(D\) 内任一路径），则 \(F(z)\) 在 \(D\) 内解析，且 \(F'(z)=f(z)\)。称 \(F\) 为 \(f\) 的\textbf{原函数}，\(F(z)+C\)（\(C\) 为复常数）为 \(f\) 的\textbf{不定积分}，记 \(\int f(z)\mathrm{d}z=F(z)+C\)。同一函数的两个原函数相差常数。

\textbf{牛顿—莱布尼茨公式}\quad 若 \(f\) 在单连通区域 \(D\) 内解析，\(F\) 是 \(f\) 的任一原函数，则对 \(z_{0},z_{1}\in D\)
\keyformula{\int_{z_{0}}^{z_{1}}f(z)\,\mathrm{d}z=F(z_{1})-F(z_{0})}
于是实积分的换元、分部积分等技巧在复积分中照旧可用（只要原函数在区域内单值解析）。

\textbf{推论}\quad 若 \(D\) 是不含零点的单连通区域，则 \(\operatorname{Ln}z\) 在 \(D\) 上存在单值解析分支，且
\[
F(z)=\ln z_{0}+\int_{z_{0}}^{z}\frac{\mathrm{d}\xi}{\xi}
\]
就是其中一个分支。更一般地，若 \(f\) 在单连通区域 \(D\) 内解析且处处不为零，则 \(\operatorname{Ln}f(z)\) 在 \(D\) 内存在单值解析分支（取 \(F'=f'/f\)）。

\subsection{复合闭路定理}
设多连通区域 \(D\) 的边界由外边界 \(C_{0}\) 与内边界 \(C_{1},C_{2},\cdots,C_{n}\) 组成，均取\textbf{正向}（沿边界前进时区域在左侧，故外边界逆时针、内边界顺时针），记复合闭路 \(C=C_{0}+C_{1}^{-}+\cdots+C_{n}^{-}\)。若 \(f\) 在 \(D\) 上解析，则
\keyformula{\oint_{C}f(z)\,\mathrm{d}z=0,\qquad\text{即}\qquad
\oint_{C_{0}}f\,\mathrm{d}z=\sum_{k=1}^{n}\oint_{C_{k}}f\,\mathrm{d}z}
（右端各内边界均取逆时针方向。）直观地说：绕外边界一周的积分等于绕各个“洞”一周的积分之和。

\corepoint{用复合闭路定理解题时\textbf{三个位置必须交代清楚}：(1) 被积函数的奇点在哪；(2) 积分曲线围住了哪些奇点；(3) 挖出的小圆周方向是逆时针还是顺时针。教材写法是“在 \(C\) 内作互不相交的小圆周 \(C_{1},\dots,C_{n}\) 分别包围各奇点”，此时 \(\oint_{C}=\sum\oint_{C_{k}}\)（都取逆时针）；若把小圆当作区域边界，则内边界是\textbf{顺时针}，符号要反过来。方向弄反是这一章最常见的错误。}

\section{柯西积分公式}

\subsection{柯西积分公式}
\textbf{定理（柯西积分公式）}\quad 设 \(f(z)\) 在简单闭曲线 \(C\) 上及其内部 \(D\) 内解析，\(z_{0}\) 为 \(D\) 内任一点，则
\keyformula{f(z_{0})=\frac{1}{2\pi\mathrm{i}}\oint_{C}\frac{f(z)}{z-z_{0}}\,\mathrm{d}z}
等价写法：\(\displaystyle\oint_{C}\frac{f(z)}{z-z_{0}}\mathrm{d}z=2\pi\mathrm{i}f(z_{0})\)。

证明要点：以 \(z_{0}\) 为心作充分小的圆 \(K:|z-z_{0}|=\rho\) 含于 \(D\)，由复合闭路定理把沿 \(C\) 的积分化为沿 \(K\) 的积分，再用 \(f\) 的连续性与估值不等式证明
\(\oint_{K}\dfrac{f(z)}{z-z_{0}}\mathrm{d}z\to 2\pi\mathrm{i}f(z_{0})\)（\(\rho\to 0\)）。

\textbf{意义}\quad 解析函数在区域\textbf{内部}的值完全由它在\textbf{边界}上的值决定。由此立得：若两个解析函数在区域边界上处处相等，则它们在整个区域上恒等（解析函数的“刚性”，实函数没有此性质）。

公式对复合闭路（多连通区域）同样成立：只要 \(z_{0}\) 在 \(C\) 所围区域内部，而 \(f\) 在 \(C\) 上及其所围区域内解析即可。

\subsection{高阶导数公式}
\textbf{定理}\quad 若 \(f(z)\) 在简单闭曲线 \(C\) 上及其内部 \(D\) 内解析，则 \(f\) 在 \(D\) 内有\textbf{任意阶导数}，且对 \(D\) 内任一点 \(z_{0}\)
\keyformula{f^{(n)}(z_{0})=\frac{n!}{2\pi\mathrm{i}}\oint_{C}\frac{f(z)}{(z-z_{0})^{n+1}}\,\mathrm{d}z,
\qquad n=1,2,\cdots}
证明思路：对 \(f(z_{0})=\frac{1}{2\pi\mathrm{i}}\oint_{C}\frac{f(z)}{z-z_{0}}\mathrm{d}z\) 关于 \(z_{0}\) 求导，用估值不等式证明“积分号下可以求导”（差商与导数的差为 \(O(|\Delta z|)\)），再对 \(n\) 归纳。

反用此式是计算积分的常用手段：
\[
\oint_{C}\frac{f(z)}{(z-z_{0})^{n+1}}\,\mathrm{d}z=\frac{2\pi\mathrm{i}}{n!}f^{(n)}(z_{0}).
\]

\corepoint{这一公式带来复分析与实分析最本质的差别：\textbf{解析 \(\Rightarrow\) 无穷次可导}（而实函数一次可导推不出二次可导）。使用时注意三点：(1) 分子 \(f\) 必须在 \(C\) 及内部\textbf{解析}——奇点要留在分母外面；(2) 分母的幂次是 \(n+1\)，对应 \(n!\)；(3) 若 \(C\) 内有多个奇点，先挖洞拆分，使每个积分只含一个奇点再套公式。}

\subsection{几个重要的推论}
\subsubsection*{（1）平均值公式}
设 \(f\) 在以 \(z_{0}\) 为心、\(R\) 为半径的圆周 \(C\) 上及其内部解析，则
\keyformula{f(z_{0})=\frac{1}{2\pi}\int_{0}^{2\pi}f\bigl(z_{0}+R\mathrm{e}^{\mathrm{i}\theta}\bigr)\,\mathrm{d}\theta
=\frac{1}{2\pi R}\oint_{C}f(z)\,\mathrm{d}s}
即解析函数在圆心的值等于它在圆周上值的平均。

\subsubsection*{（2）柯西不等式}
设 \(M=\max\limits_{z\in C}|f(z)|\)，则
\keyformula{\left|f^{(n)}(z_{0})\right|\leqslant\frac{n!\,M}{R^{n}};\qquad
\text{特别地}\ \left|f'(z_{0})\right|\leqslant\frac{M}{R}}
它把函数在一点处的各阶导数用它在圆周上的最大模控制住，是估计导数、证明刘维尔定理的工具。

\subsubsection*{（3）刘维尔（Liouville）定理}
\textbf{整函数}指在复平面上处处解析的函数（多项式、\(\mathrm{e}^{z}\)、\(\sin z\) 等）。若整函数 \(f(z)\) 在复平面上\textbf{有界}（\(|f(z)|\leqslant M\)），则 \(f\) 必为常数。

证明：对任意 \(z\) 与任意 \(R>0\)，由柯西不等式 \(|f'(z)|\leqslant M/R\)；令 \(R\to+\infty\) 得 \(f'(z)=0\)，故 \(f\) 为常数。

\subsubsection*{（4）代数基本定理}
任意 \(n\) 次（\(n\geqslant 1\)）复系数多项式
\[
f(z)=a_{0}z^{n}+a_{1}z^{n-1}+\cdots+a_{n-1}z+a_{n}\qquad (a_{0}\neq 0)
\]
在复平面上\textbf{至少有一个零点}，从而恰有 \(n\) 个零点（计重数）。

证明：若 \(f\) 无零点，则 \(F=1/f\) 是整函数，且由 \(|f(z)|\to+\infty\ (z\to\infty)\) 知 \(F\) 有界，由刘维尔定理 \(F\) 为常数，矛盾。

\subsubsection*{（5）莫累拉（Morera）定理}
设 \(f(z)\) 在区域 \(D\) 内连续，且对 \(D\) 内任一条其内部含于 \(D\) 的简单闭曲线 \(C\) 都有 \(\oint_{C}f\,\mathrm{d}z=0\)，则 \(f\) 在 \(D\) 内解析。

它是柯西积分定理的\textbf{逆定理}：柯西定理用“解析”推出“积分为零”，莫累拉定理用“积分为零”反推“解析”，常用于证明由积分或级数定义的函数是解析的。

\subsubsection*{（6）泊松（Poisson）积分公式}
设 \(f\) 在圆 \(C:|z-z_{0}|=R\) 内及 \(C\) 上解析，\(z=z_{0}+r\mathrm{e}^{\mathrm{i}\theta}\ (r<R)\) 为圆内一点，则
\[
f(z)=\frac{1}{2\pi}\int_{0}^{2\pi}
\frac{R^{2}-r^{2}}{R^{2}-2Rr\cos(\varphi-\theta)+r^{2}}\,
f\bigl(z_{0}+R\mathrm{e}^{\mathrm{i}\varphi}\bigr)\,\mathrm{d}\varphi .
\]
比较两端的实部即得\textbf{泊松积分公式}：圆内调和函数的值可由它在圆周上的值通过上述积分表示。特别取 \(r=0\) 即回到调和函数的平均值公式。

\section{解析函数模的性质}

下面的结论由平均值公式与柯西不等式得到（讲义 §1.4），它刻画了解析函数模的“边界控制”性质。

\textbf{最大模原理}\quad 设 \(D\) 是有界连通区域，\(f(z)\) 在 \(\overline{D}\) 上连续、在 \(D\) 内解析，则
\keyformula{\max_{z\in\overline{D}}|f(z)|=\max_{z\in\partial D}|f(z)|=M}
并且：若存在 \(D\) 的\textbf{内点} \(z_{0}\) 使 \(|f(z_{0})|=M\)，则 \(f\) 在 \(D\) 内为常数。

证明要点：设内点取到最大值 \(M\)，由平均值公式
\(M=|f(z_{0})|\leqslant\frac{1}{2\pi}\int_{0}^{2\pi}|f(z_{0}+r\mathrm{e}^{\mathrm{i}\theta})|\mathrm{d}\theta\leqslant M\)，
故等号处处成立、\(|f|\) 在圆内恒为 \(M\)；沿折线逐点推过去，得 \(|f|\) 在 \(D\) 内恒为 \(M\)，再由 §2.2 的推论知 \(f\) 为常数。

\textbf{最小模原理}\quad 设 \(D\) 有界连通，\(f\) 在 \(\overline{D}\) 上连续、在 \(D\) 内解析且\textbf{处处不为零}，则
\[
\min_{z\in\overline{D}}|f(z)|=\min_{z\in\partial D}|f(z)|;
\]
若在内点取到最小值，则 \(f\) 为常数。（对 \(1/f\) 用最大模原理。）

\textbf{调和函数的极值原理}\quad 若 \(u(x,y)\) 是有界闭区域 \(\overline{D}\) 上的调和函数（在 \(D\) 内调和、在 \(\overline{D}\) 上连续），则 \(u\) 的最大值与最小值都在边界 \(\partial D\) 上取到；若在内点取到，则 \(u\) 为常数。

\corepoint{最大模原理管的是 \(|f|\) 而\textbf{不是} \(f\)：\(f\) 本身（复值）谈不上大小。三个使用前提别漏：区域\textbf{有界}、\(f\) 在闭区域上\textbf{连续}、在内部\textbf{解析}。典型应用：证明“模为常数的解析函数是常数”（\(|f|\equiv c\) 时内部也取到最大值）。对照两个二元实函数能看清“调和”的关键：\(u=x^{2}+y^{2}\) 在圆盘上的最大值在圆心取到（\(\Delta u=4\neq0\)，不调和，不受定理约束）；而调和的 \(u=x^{2}-y^{2}=\operatorname{Re}(z^{2})\) 的最大值必在边界上取到。}
